\section{Compact factorization and the approximation property}
\label{sec:approximation}

We now transfer Walsh decay to the identity of \(X\).
The approximation property will let us apply Theorem~\ref{thm:uniform-decay}
to compact maps into \(X\). We therefore factor the inclusion of
each finite-dimensional \(E\subset X\) through a compact map
whose domain dual has a controlled moment estimate.
Using \(E\) itself as the domain would not suffice: the quotient
\(E^*\) of \(X^*\) need not inherit the cotype estimate of \(X^*\).
Instead, we enlarge \(E\) to retain the moment lifts available in
\(X\), and use decreasing weights to obtain compactness.

\begin{lemma}[A compact domain with controlled dual moments]
\label{lem:compact-domain}
Let \(m\ge1\) be an integer, let \(\Delta\ge1\), and let \(X\) be
a Banach space with \(c_m(X^*)\le\Delta\).
For every nonzero finite-dimensional subspace \(E\subset X\), there
exist a Banach space \(G\), an injective compact operator \(J:G\to X\),
and a contraction \(\iota:E\to G\), such that
\begin{equation}\label{eq:compact-domain}
J\iota x=x\quad(x\in E),\qquad
\|J\|\le2,\qquad c_m(G^*)\le12\Delta.
\end{equation}
\end{lemma}
\begin{proof}
\textit{Finite-dimensional enlargements.}
We construct nested finite-dimensional subspaces
\[
E=E_0\subset E_1\subset E_2\subset\cdots\subset X
\]
such that each \(z=(z_i)_{i=1}^m\in\ell_2^m(E_l)\) has a lift
\(H_l:\{-1,1\}^m\to E_{l+1}\) satisfying
\begin{equation}\label{eq:finite-level-lift}
\E_\eps\eps_iH_l(\eps)=z_i,\qquad
\|H_l\|_{L_2(E_{l+1})}\le3\Delta\|z\|_{\ell_2^m(E_l)}.
\end{equation}
First take moments in the whole space \(X\). The map
\[
U:L_2(\{-1,1\}^m;X)\to\ell_2^m(X),\qquad
UH=(\E_\eps\eps_iH(\eps))_i,
\]
has the bounded right inverse
\((z_i)\mapsto\sum_i\eps_i z_i\), and
\[
\|U^*(\theta_i)\|_{L_2(X^*)}
=\left(\E_\eps\left\|\sum_i\eps_i\theta_i\right\|^2\right)^{1/2}
\ge\Delta^{-1}\|(\theta_i)\|_{\ell_2^m(X^*)}.
\]
By \eqref{eq:quotient-lift}, the infimal lift norm of a unit tuple
is at most \(\Delta\), so a lift of norm below \(2\Delta\) exists.
Take a finite \(\eta\)-net on the unit sphere of \(\ell_2^m(E_l)\),
where \(0<\eta\le\min(1/2,\Delta/\sqrt m)\), and choose such a lift
for each net point. Let \(E_{l+1}\) be the span of \(E_l\) and all
the values of these finitely many lifts.

For any unit tuple \(z\in\ell_2^m(E_l)\), choose a net point \(v\)
with \(\|z-v\|_{\ell_2^m(X)}\le\eta\). Add
\(\sum_i\eps_i(z_i-v_i)\) to its lift. This corrects the moments
exactly and adds at most
\[
\left\|\sum_i\eps_i(z_i-v_i)\right\|_{L_2(X)}
\le\sqrt m\,\|z-v\|_{\ell_2^m(X)}\le\Delta
\]
to the norm. This proves \eqref{eq:finite-level-lift} for unit
tuples; scaling and the zero lift cover the other cases.

\medskip
\textit{The compact map.}
Put
\[
\mathcal H=\left(\bigoplus_{l=0}^\infty E_l\right)_{\ell_2},
\qquad
j:\mathcal H\to X,\qquad j(z)=\sum_{l\ge0}2^{-l}z_l.
\]
Cauchy--Schwarz proves absolute convergence of the series and gives
\[
\|j\|\le\left(\sum_{l\ge0}4^{-l}\right)^{1/2}<2.
\]
The same estimate on the tails shows that the finite-rank truncations
converge in operator norm. Thus \(j\) is compact.
Let \(G=\mathcal H/\ker j\), with quotient map \(Q\), and define
\(J\) by \(JQ=j\). The kernel is closed, so \(G\) is Banach;
the quotient norm gives \(\|J\|\le\|j\|<2\), and \(J\) is injective.
Every element of \(B_G\) has a representative in \(2B_{\mathcal H}\).
Hence \(J(B_G)\subset j(2B_{\mathcal H})\) is relatively compact,
which proves compactness of \(J\). Finally,
\(\iota(x)=Q(x,0,0,\ldots)\) is a contraction and \(J\iota x=x\).
The finite-rank truncations of \(j\) need not annihilate \(\ker j\),
so they need not induce operators on \(G\). Thus this compactness
argument does not itself approximate \(J\) by finite-rank maps;
that passage will use the approximation property of \(X\).

\medskip
\textit{Dual moments from the weighted shift.}
The nesting of the \(E_l\)'s makes
\[
R:\mathcal H\to\mathcal H,\qquad
R(z_0,z_1,\ldots)=(0,2z_0,2z_1,\ldots)
\]
a bounded map. The factor \(2\) compensates for the next dyadic
weight, so \(jR=j\) and therefore \(QR=Q\).
This identity lets us assemble the finite-level lifts into exact
moment lifts in \(G\).

For \(g_1,\ldots,g_m\in G\), choose representatives
\(z^i=(z_l^i)_l\in\mathcal H\) with
\(Qz^i=g_i\) and \(\|z^i\|\le2\|g_i\|\), taking \(z^i=0\)
when \(g_i=0\).
Use \eqref{eq:finite-level-lift} on each tuple \((z_l^i)_i\),
and set
\[
\widetilde H(\eps)=(0,2H_0(\eps),2H_1(\eps),\ldots).
\]
Summing squared estimates yields
\[
\|\widetilde H\|_{L_2(\mathcal H)}
\le6\Delta\left(\sum_i\|z^i\|^2\right)^{1/2}
\le12\Delta\left(\sum_i\|g_i\|^2\right)^{1/2}.
\]
Every atom of the sign cube has positive mass, so this estimate
also shows that \(\widetilde H(\eps)\in\mathcal H\) for every
\(\eps\). Its \(i\)-th moment is \(Rz^i\). Since \(QR=Q\),
\[
\E_\eps\eps_iQ\widetilde H(\eps)=QRz^i=Qz^i=g_i.
\]
For \(\theta_1,\ldots,\theta_m\in G^*\), it follows that
\begin{align*}
\left|\sum_i\theta_i(g_i)\right|
&=\left|\E_\eps
\left\langle Q\widetilde H(\eps),\sum_i\eps_i\theta_i\right\rangle\right|\\
&\le12\Delta
\left(\sum_i\|g_i\|^2\right)^{1/2}
\left(\E_\eps\left\|\sum_i\eps_i\theta_i\right\|^2\right)^{1/2}.
\end{align*}
Taking the supremum over \(\sum_i\|g_i\|^2\le1\) gives
\(c_m(G^*)\le12\Delta\) by finite \(\ell_2\) duality.
\end{proof}

We now complete the proof of the main theorem.

\begin{proof}[Proof of Theorem~\ref{thm:main}]
First suppose \(K(X)<\infty\).
On each finite cube, the adjoint of the Rademacher projection on
\(L_2(X)\) is the corresponding projection on \(L_2(X^*)\).
Therefore \(K(X^*)=K(X)\).
By the classical characterization recalled in the introduction,
both \(X\) and \(X^*\) have type strictly greater than one.

If a space \(Y\) has type \(p>1\), with constant \(D\), then \(Y^*\)
has cotype \(p'=p/(p-1)\).
Indeed, for finite tuples \(y_i\in Y\), \(y_i^*\in Y^*\),
\begin{align*}
\left|\sum_i y_i^*(y_i)\right|
&=\left|\E_\eps
\left\langle\sum_i\eps_i y_i,\sum_i\eps_i y_i^*\right\rangle\right|\\
&\le D\left(\sum_i\|y_i\|^p\right)^{1/p}
\left(\E_\eps\left\|\sum_i\eps_i y_i^*\right\|^2\right)^{1/2}.
\end{align*}
Taking the supremum over the unit ball of \(\ell_p^m(Y)\) proves
the claim. Apply it to \(Y=X\) and \(Y=X^*\), and restrict the
resulting cotype inequality on \(X^{**}\) to \(X\).
This direction does not use the approximation property.

Conversely, suppose that \(X\) has the approximation property and
\[
C_q(X)\le C_F,\qquad C_r(X^*)\le C_Y
\]
for finite \(q,r\ge2\) and finite constants \(C_F,C_Y\ge1\).
For an integer \(m\), set
\[
D_f=C_Fm^{1/2-1/q},\qquad
\Delta=C_Ym^{1/2-1/r},\qquad D_g=12\Delta.
\]
The parameters in \eqref{eq:lift-parameters} then satisfy
\[
ma=\frac{m^{1/r}}{24LC_Y},\qquad
mb=\frac{m^{1/q}}{2LC_F}.
\]
These positive powers of \(m\) eventually exceed any fixed multiple
of \((\log m)^3\). Choose \(m\) once so that the hypotheses of
Lemma~\ref{lem:coefficient} hold, and fix the decay sequence \(\rho_N\)
given by Theorem~\ref{thm:uniform-decay} for \(m,D_g,D_f\).
All these parameters are independent of the finite-dimensional
subspaces and approximating operators chosen below.

We first show that
\begin{equation}\label{eq:identity-walsh-decay}
w_N(\Id_X)\le2\rho_N\qquad(N\ge0).
\end{equation}
For a nonzero finite-dimensional \(E\subset X\), choose
\(J:G\to X\) and \(\iota:E\to G\) by Lemma~\ref{lem:compact-domain}.
Apply the approximation property to the compact set
\(M=\overline{J(B_G)}\). For every \(j\ge1\), it supplies a bounded
finite-rank map \(S_j:X\to X\) such that
\[
\|S_jJ-J\|
\le\sup_{x\in M}\|S_jx-x\|<1/j.
\]
Thus \(J\) is an operator-norm limit of finite-rank maps from
\(G\) to \(X\). Since \(c_m(G^*)\le D_g\) and \(c_m(X)\le D_f\),
the norm-limit assertion of Theorem~\ref{thm:uniform-decay} gives
\begin{equation}\label{eq:compact-walsh-decay}
w_N(J)\le\rho_N\|J\|\le2\rho_N.
\end{equation}
Only the composites \(S_jJ\) converge in operator norm;
no common bound on \(\|S_j\|\) is required.

Now fix any input \(x=(x_u)_{u\in V_N}\in L_2(V_N;X)\).
If \(x=0\), there is nothing to prove. Otherwise take
\(E=\Span\{x_u:u\in V_N\}\) and choose the corresponding \(J,\iota\).
The family \(z_u=\iota(x_u)\) satisfies
\(\|z\|_{L_2(G)}\le\|x\|_{L_2(X)}\) and \(Jz_u=x_u\).
Consequently,
\[
\|W_N^{\Id_X}x\|_{L_2(X)}
=\|W_N^Jz\|_{L_2(X)}
\le2\rho_N\|x\|_{L_2(X)},
\]
which proves \eqref{eq:identity-walsh-decay}.

If \(K(X)=\infty\), the characterization
\eqref{eq:classical-characterization} supplies \(\lambda\ge1\)
such that, for every \(N\ge1\), there are \(x_u\in B_X\), \(u\in V_N\),
with
\[
\left\|\sum_{u\in V_N}\alpha_u x_u\right\|
\ge\lambda^{-1}\sum_{u\in V_N}|\alpha_u|
\]
for all real coefficients. These vectors form an input of \(L_2\)
norm at most one, while every Walsh output satisfies
\[
\|(W_N^{\Id_X}x)(v)\|
=\left\|2^{-N}\sum_{u\in V_N}\chi(u,v)x_u\right\|
\ge\lambda^{-1}.
\]
Thus \(w_N(\Id_X)\ge\lambda^{-1}\) for every \(N\), contradicting
\eqref{eq:identity-walsh-decay} and \(\rho_N\to0\).
This proves \(K(X)<\infty\).
\end{proof}

\begin{proof}[Proof of Corollary~\ref{cor:uniform-cotype-entropy}]
If no uniform \(\kappa\) existed, choose spaces \(X_j\) with the
stated fixed cotype bounds and AP, but \(K(X_j)\to\infty\), and set
\[
Z=\left(\bigoplus_{j\ge1}X_j\right)_{\ell_2}.
\]
Finite-coordinate truncations converge uniformly on each compact subset
of \(Z\). Approximating its finitely many compact coordinate images by
finite-rank maps therefore proves AP for \(Z\); no uniform bound on
the norms of those maps is needed. Moreover,
\(Z^*=(\bigoplus_{j\ge1}X_j^*)_{\ell_2}\) isometrically.
For any finite tuple \(z_i=(z_{ij})_j\) in \(Z\), Minkowski's inequality
in \(\ell_{q/2}\) and the coordinate cotype bounds give
\[
\left(\sum_i\|z_i\|_Z^q\right)^{2/q}
\le \sum_j\left(\sum_i\|z_{ij}\|^q\right)^{2/q}
\le C^2\E_\eps\left\|\sum_i\eps_i z_i\right\|_Z^2.
\]
Thus \(C_q(Z)\le C\), and the same argument on the dual gives
\(C_r(Z^*)\le D\). Theorem~\ref{thm:main} implies \(K(Z)<\infty\).
Restriction of each finite-cube Rademacher projection to a coordinate
summand gives \(K(X_j)\le K(Z)\), a contradiction.

For the entropy assertion, finite-dimensional gauge spaces have AP,
so the first assertion bounds their \(K\)-constants by \(\kappa\).
The bounded-\(K\) entropy-duality theorem of
Artstein--Milman--Szarek--Tomczak-Jaegermann \cite[Theorem~5]{AMSTJ2004},
in the covering-number formulation of
\cite[Theorem~1.1, equation~(1.6), and Section~2]{Milman2007},
has constants \(a,b\ge1\) depending only on an upper bound for
the \(K\)-convexity constant of either gauge space.
Polarity identifies \(X_{B^\circ}=X_B^*\), and the finite-cube
adjoint identity used above gives \(K(X_B^*)=K(X_B)\).
Thus the same bound is available for the polar bodies in the reverse
inequality. Also, \(X_{tL}\) is linearly isometric to \(X_L\).
Apply the covering theorem to \(K,tL\). Since
\((tL)^\circ=t^{-1}L^\circ\) and
\(N(t^{-1}L^\circ,cK^\circ)=N(L^\circ,ctK^\circ)\) for \(c>0\),
both inequalities in \eqref{eq:cotype-entropy-duality} follow for every
\(t>0\), with constants depending only on \(\kappa\).
\end{proof}
